\documentclass[11pt]{amsart}
\usepackage[localtheorems]{raeez-math-template}

\providecommand{\C}{\mathbb{C}}
\providecommand{\bP}{\mathbb{P}}
\providecommand{\bY}{\mathbb{Y}}
\providecommand{\cO}{\mathcal{O}}
\providecommand{\Rep}{\mathrm{Rep}}
\providecommand{\Tot}{\mathrm{Tot}}
\providecommand{\Eone}{E_1}
\providecommand{\Etwo}{E_2}
\providecommand{\CoHA}{\mathrm{CoHA}}
\providecommand{\DT}{\mathrm{DT}}
\providecommand{\tr}{\mathrm{tr}}
\providecommand{\kch}{\kappa_{\mathrm{ch}}}
\providecommand{\kcat}{\kappa_{\mathrm{cat}}}
\providecommand{\kBKM}{\kappa_{\mathrm{BKM}}}
\providecommand{\kfib}{\kappa_{\mathrm{fiber}}}

\theoremstyle{plain}
\newtheorem{theorem}{Theorem}[section]
\newtheorem{proposition}[theorem]{Proposition}
\newtheorem{lemma}[theorem]{Lemma}
\newtheorem{corollary}[theorem]{Corollary}
\newtheorem{conjecture}[theorem]{Conjecture}
\theoremstyle{definition}
\newtheorem{definition}[theorem]{Definition}
\newtheorem{construction}[theorem]{Construction}
\theoremstyle{remark}
\newtheorem{remark}[theorem]{Remark}

\title{The conifold positive half and its chiral realisation}
\date{}

\begin{document}
\maketitle

The resolved conifold
\[
  \bY=\Tot\bigl(\cO_{\bP^1}(-1)\oplus\cO_{\bP^1}(-1)\bigr)
\]
is a smooth quasi-projective toric Calabi-Yau threefold with one compact
curve.  The Hall object is the unframed critical CoHA of the two-vertex
Klebanov-Witten Jacobi algebra.  It is an \(\Eone\)
algebra \(A_{\mathrm{con}}\).  Its quasitriangular and vertex-algebra
data live on \(D(A_{\mathrm{con}})\), on
\(\mathcal{Z}(\Rep^{\Eone}(A_{\mathrm{con}}))\), or on evaluation
modules of the double.

\section{The conifold NCCR}

\begin{construction}[Holomorphic volume form]\label{cons:conifold-volume-form}
Let \(U_0=\C_t\times\C^2_{u,v}\) and
\(U_\infty=\C_{\tilde t}\times\C^2_{\tilde u,\tilde v}\) be the standard
cover of \(\bY\), with
\[
  \tilde t=t^{-1},\qquad \tilde u=t u,\qquad \tilde v=t v .
\]
Then
\[
  dt\wedge du\wedge dv
  =
  -\,d\tilde t\wedge d\tilde u\wedge d\tilde v
\]
on \(U_0\cap U_\infty\).  Hence
\[
  \Omega_{\bY}|_{U_0}=dt\wedge du\wedge dv,\qquad
  \Omega_{\bY}|_{U_\infty}=-\,d\tilde t\wedge d\tilde u\wedge d\tilde v
\]
defines a nowhere-vanishing holomorphic \(3\)-form.  Equivalently,
\[
  c_1(T\bY)=c_1(T\bP^1)+c_1(\cO(-1))+c_1(\cO(-1))=2-1-1=0 .
\]
\end{construction}

\begin{construction}[Klebanov-Witten Jacobi algebra]\label{cons:kw-jacobi}
Let \(Q_{\mathrm{KW}}\) have vertices \(0,1\), arrows
\[
  a_1,a_2\colon 0\to 1,\qquad b_1,b_2\colon 1\to 0,
\]
and potential
\[
  W_{\mathrm{KW}}
  =
  \tr(a_1b_1a_2b_2-a_1b_2a_2b_1).
\]
The Jacobi algebra
\[
  J_{\mathrm{KW}}
  =
  \C Q_{\mathrm{KW}}\big/
  \langle \partial_a W_{\mathrm{KW}}:a\in Q_{\mathrm{KW},1}\rangle
\]
is the Van den Bergh-Szendr\H{o}i noncommutative crepant resolution of
the affine conifold \(X_0=\Spec \C[x,y,z,w]/(xy-zw)\).  The four cyclic
derivatives are
\[
\begin{aligned}
  \partial_{a_1}W_{\mathrm{KW}}&=b_1a_2b_2-b_2a_2b_1,\\
  \partial_{a_2}W_{\mathrm{KW}}&=b_2a_1b_1-b_1a_1b_2,\\
  \partial_{b_1}W_{\mathrm{KW}}&=a_2b_2a_1-a_1b_2a_2,\\
  \partial_{b_2}W_{\mathrm{KW}}&=a_1b_1a_2-a_2b_1a_1 .
\end{aligned}
\]
On the vertex \(0\) corner one normalization sends the mesons
\[
  x=a_1b_1,\qquad y=a_2b_2,\qquad z=a_1b_2,\qquad w=a_2b_1
\]
to the four affine coordinates.  The corresponding vertex \(1\) mesons
represent the same central classes by the Jacobi relations, and
\[
  Z(J_{\mathrm{KW}})=\C[x,y,z,w]/(xy-zw).
\]
The two small resolutions of \(X_0\) are the two geometric chambers of
this NCCR.
\end{construction}

\section{The positive half}

\begin{definition}[Conifold critical CoHA]\label{def:conifold-coha}
For \(\gamma=(\gamma_0,\gamma_1)\in\mathbb{N}^2\), set
\[
  \Rep_\gamma(Q_{\mathrm{KW}})
  =
  \bigoplus_{a\in Q_{\mathrm{KW},1}}
  \mathrm{Hom}\bigl(\C^{\gamma_{s(a)}},\C^{\gamma_{t(a)}}\bigr),
  \qquad
  G_\gamma=\mathrm{GL}_{\gamma_0}\times\mathrm{GL}_{\gamma_1}.
\]
The trace of \(W_{\mathrm{KW}}\) defines a \(G_\gamma\)-invariant
function \(\tr(W_{\mathrm{KW}})_\gamma\) on \(\Rep_\gamma(Q_{\mathrm{KW}})\).
The critical cohomological Hall algebra is
\[
  \CoHA_{\mathrm{con}}
  =
  \bigoplus_{\gamma\in\mathbb{N}^2}
  H^{*}_{c,G_\gamma}
  \bigl(\Rep_\gamma(Q_{\mathrm{KW}}),
  \varphi_{\tr(W_{\mathrm{KW}})_\gamma}\bigr)^\vee ,
\]
with product given by the Kontsevich-Soibelman correspondence.
\end{definition}

\begin{construction}[Two-colour shuffle kernel]\label{cons:conifold-shuffle-kernel}
Let \(h_1,h_2\) be the torus weights on the two fibre directions and set
\(h_3=-h_1-h_2\).  In additive shuffle variables \(u=z-w\), the
Klebanov-Witten cross-kernel is
\[
  \omega_{01}(u)=
  \frac{(u+h_1)(u+h_2)}
       {u(u+h_1+h_2)},
  \qquad
  \omega_{10}(u)=
  \frac{(u-h_1)(u-h_2)}
       {u(u-h_1-h_2)} .
\]
The diagonal kernels are \(1\).  The two numerator factors record the
two arrows in a fixed direction.  The poles at \(u=0\) and
\(u=\mp(h_1+h_2)\) record the diagonal morphism and the
CY\(_3\)-Serre-dual Jacobi obstruction.  This is the rank-two
conifold kernel, distinct from the one-colour \(\C^3\) cubic kernel.
\end{construction}

\begin{theorem}[Conifold positive half]\label{thm:conifold-positive-half}
After the standard cohomological shift and shuffle completion, the
Hall-side conifold algebra is the positive half
\[
  \CoHA_{\mathrm{con}}
  \simeq
  Y^+\!\bigl(\widehat{\mathfrak{gl}}(1|1)\bigr)^{\mathrm{con}}
\]
as a \(\mathbb{N}^2\)-graded connected bialgebra in super-vector spaces
over \(\C((\hbar))\), with positive-root set
\[
  \Delta_+^{\mathrm{re}}
  =
  \{(n+1,n),(n,n+1):n\geq 0\},\qquad
  \Delta_+^{\mathrm{im}}=\{(n,n):n\geq 1\}.
\]
The real root sectors have primitive supertrace \(1\); the imaginary
sector has primitive supertrace \(-2\).  A nondegenerate Hall pairing
between the positive and negative halves gives the Drinfeld double
\[
  D(\CoHA_{\mathrm{con}})
  \simeq
  Y\!\bigl(\widehat{\mathfrak{gl}}(1|1)\bigr)^{\mathrm{con}},
\]
and the braided \(\Etwo\) structure is the corresponding centre or
double layer.
\end{theorem}

\begin{proof}
The Jacobi algebra in Construction~\ref{cons:kw-jacobi} gives the
conifold NCCR.  Davison's critical CoHA construction applies to the
quiver with potential \((Q_{\mathrm{KW}},W_{\mathrm{KW}})\), and the
two-colour kernel in Construction~\ref{cons:conifold-shuffle-kernel}
is Negu\c{t}'s conifold shuffle kernel.  Davison-Meinhardt PBW
identifies the primitive mixed Hodge structure with the conifold BPS
Lie algebra.  Morrison-Mozgovoy-Nagao-Szendr\H{o}i compute the chamber
decomposition and the affine \(A_1^{(1)}\) root system:
\[
  \Omega(n+1,n)=\Omega(n,n+1)=1,\qquad
  \Omega(n,n)=-2 .
\]
Li-Yamazaki give the corresponding current presentation as the
positive half of the affine super Yangian of
\(\widehat{\mathfrak{gl}}(1|1)\).  The Hall product is associative and
ordered, hence \(\Eone\).  A universal \(R\)-matrix requires the Hall
pairing and the negative half, i.e. the Drinfeld double, or equivalently
the centre of the \(\Eone\)-monoidal representation category.
\end{proof}

\begin{proposition}[Chart algebra and global two-colour algebra]\label{prop:chart-restriction}
Each toric chart \(U_\pm\simeq\C^3\) carries the one-vertex critical
CoHA
\[
  \CoHA(\C^3)\simeq Y^+(\widehat{\mathfrak{gl}}_1)
\]
of Schiffmann-Vasserot and Tsymbaliuk.  The global conifold algebra has
three inputs: the two chart algebras
\[
  \CoHA(U_+)\simeq Y^+(\widehat{\mathfrak{gl}}_1),
  \qquad
  \CoHA(U_-)\simeq Y^+(\widehat{\mathfrak{gl}}_1),
\]
the Laurent transition on \(U_+\cap U_-\simeq\C^*\times\C^2\),
\[
  (t,u,v)\longmapsto (t^{-1},tu,tv),
\]
and the cross-kernel of Construction~\ref{cons:conifold-shuffle-kernel}.
The compact-curve sector \(\Delta_+^{\mathrm{im}}=\{(n,n):n\geq 1\}\)
belongs to this global two-colour algebra.
\end{proposition}

\begin{proof}
The fan of \(\bY\) has two maximal cones, each spanned by a lattice
basis; the corresponding affine toric varieties are \(\C^3\).  The
critical CoHA of the one-vertex three-loop quiver with potential
\(\tr(x[y,z])\) is \(Y^+(\widehat{\mathfrak{gl}}_1)\).  The common wall
of the two cones gives the overlap \(\C^*\times\C^2\).  The transition
matrix has the determinant cancellation computed in
Construction~\ref{cons:conifold-volume-form}; the remaining data are
the Fourier-Mukai twist and the two-colour cross-kernel.  These data
produce the \((n,n)\) compact-curve tower in
Theorem~\ref{thm:conifold-positive-half}.
\end{proof}

\section{Chiral layer}

\begin{proposition}[Operadic level]\label{prop:conifold-operadic-level}
For the conifold CY\(_3\) input, the native output of the
CY-to-chiral construction is \(\Eone\)-chiral.  Any \(\Etwo\) structure
attached to it belongs to the representation-theoretic layer
\[
  \mathcal{Z}\bigl(\Rep^{\Eone}(A_{\mathrm{con}})\bigr)
  \quad\text{or}\quad
  D(\CoHA_{\mathrm{con}}),
\]
with the half-braiding encoded by the Yangian \(R\)-matrix.
\end{proposition}

\begin{proof}
The dimension rule in the CY-to-chiral construction gives native
operadic level
\[
  n_{\mathrm{native}}(d)=
  \begin{cases}
  \infty, & d=1,\\
  2, & d=2,\\
  1, & d\geq 3 .
  \end{cases}
\]
Since \(\bY\) has complex dimension \(3\), its native chiral output is
\(\Eone\).  The \(\Etwo\) layer at \(d\geq 3\) is the Drinfeld-centre
layer of the representation category, in agreement with the double in
Theorem~\ref{thm:conifold-positive-half}.
\end{proof}

\begin{conjecture}[Holomorphic Chern-Simons realisation]\label{conj:hcs-realisation}
Abelian holomorphic Chern-Simons theory on \(\bY\), restricted to a
formal neighbourhood of the exceptional curve \(C=\bP^1\), admits a
comparison map to a factorisation algebra whose Hall shadow is
\(\CoHA_{\mathrm{con}}\).  On the two affine charts this comparison
recovers the \(\C^3\) affine-Yangian structure function
\[
  \varphi_{\C^3}(z)=
  \frac{(z+h_1)(z+h_2)(z+h_3)}
       {(z-h_1)(z-h_2)(z-h_3)},
  \qquad h_1+h_2+h_3=0 .
\]
The global conifold realisation is obtained by descent across
\((t,u,v)\mapsto(t^{-1},tu,tv)\), with the compact curve producing the
imaginary-root sector of \(Y^+(\widehat{\mathfrak{gl}}(1|1))^{\mathrm{con}}\).
\end{conjecture}

\begin{remark}[Proof boundary for the hCS comparison]\label{rem:hcs-proof-boundary}
The proved input is the Klebanov-Witten critical CoHA and the chart
\(\C^3\) affine-Yangian calculation.  The conjectural input is the
oriented comparison from perturbative holomorphic Chern-Simons
observables on the formal neighbourhood of \(C\) to the critical
vanishing-cycle CoHA.  The comparison must preserve the Laurent
transition, the two-colour kernel, and the Hall product.
\end{remark}

\begin{remark}[Layer separation]\label{rem:no-w-confusion}
The equality \(\CoHA(\C^3)=Y^+(\widehat{\mathfrak{gl}}_1)\) is a
chart statement about the positive half.  A vertex algebra of
\(\mathcal{W}_{1+\infty}\)-type requires a double, centre, completion,
or evaluation representation.  The conifold chain is
\[
  Y^+\!\bigl(\widehat{\mathfrak{gl}}(1|1)\bigr)^{\mathrm{con}}
  \hookrightarrow
  D\!\left(Y^+\!\bigl(\widehat{\mathfrak{gl}}(1|1)\bigr)^{\mathrm{con}}\right)
  \xrightarrow{\mathrm{ev}_\lambda}
  \End(\mathcal V_\lambda),
\]
where \(\mathcal V_\lambda\) is the chosen free-field or
\(\mathcal{W}\)-type vacuum module.  The Hall algebra is the first
term of this chain.
\end{remark}

\section{The chiral invariant}

\begin{corollary}[Direct conifold value]\label{cor:conifold-kappach}
The chiral invariant of the conifold Hall shadow is
\[
  \kch(A_{\mathrm{con}})=\DT_{(1,0)}=1 .
\]
\end{corollary}

\begin{proof}
The dimension vector \((1,0)\) is a simple real root of the
Klebanov-Witten quiver.  At this dimension all arrows vanish and
\(\tr(W_{\mathrm{KW}})\) is zero.  The primitive BPS summand is the
simple module at vertex \(0\); after the Davison cohomological shift and
PBW extraction it contributes one generator.  The
Morrison-Mozgovoy-Nagao-Szendr\H{o}i root decomposition gives
\(\Omega(1,0)=1\).  This is the direct McKay/Jacobi/NCCR computation of
the conifold Hall shadow.  In the resolved chamber it is the genus-zero
BPS multiplicity of the single primitive compact curve class
\([\bP^1]\).
\end{proof}

\begin{remark}[Separation of the four invariants]\label{rem:four-invariants}
The value in Corollary~\ref{cor:conifold-kappach} is
\(\kch\).  The compact categorical invariant
\(\kcat=\chi(\cO_X)\) lies outside its compact-CY definition on the
quasi-projective total space \(\bY\), since
\(H^0(\bY,\cO_{\bY})\) is infinite-dimensional.  The Borcherds invariant
\(\kBKM=c_N(0)/2\) requires a Borcherds product input.  The fibre
invariant \(\kfib\) requires a chosen compact fibre or lattice
datum.  The conifold Jacobi/NCCR computation supplies exactly
\(\kch(A_{\mathrm{con}})=1\).
\end{remark}

\begin{remark}[Surface canonical-bundle formula]\label{rem:surface-canonical-bundle-formula}
The surface canonical-bundle formula
\[
  \kappa_{\mathrm{ch}}\bigl(\Tot(\omega_S)\bigr)
  =
  \frac{1}{2}\chi_{\mathrm{top}}(S)
\]
has input a compact Fano surface \(S\).  The conifold has base
\(\bP^1\) and rank-two normal bundle \(\cO(-1)^{\oplus 2}\).  Thus the
value \(3/2\) belongs to the local \(\bP^2\) computation alone.
\end{remark}

\paragraph{Primary anchors.}
Van den Bergh 2004, \emph{Three-dimensional flops and non-commutative
rings}, Theorem 3.2.10; Szendr\H{o}i 2008, \emph{Non-commutative
Donaldson-Thomas theory and the conifold}; Morrison, Mozgovoy, Nagao,
and Szendr\H{o}i 2011, \emph{Motivic Donaldson-Thomas invariants of the
conifold and the refined topological vertex}, Theorem 1; Davison
2013, \texttt{arXiv:1311.6989}, Theorem A; Davison-Meinhardt 2016,
\texttt{arXiv:1601.02479}, Theorem A; Negu\c{t} 2015,
\texttt{arXiv:1512.06473}, equation (1.6); Li-Yamazaki 2020,
\texttt{arXiv:2003.08909}, Section 8.3.6.3; Schiffmann-Vasserot 2013,
\emph{Cherednik algebras, W-algebras and the equivariant cohomology of
the moduli space of instantons on \(\mathbb{A}^2\)}; Tsymbaliuk 2017,
\emph{The affine Yangian of \(\mathfrak{gl}_1\) revisited}; Costello
2013, \emph{Supersymmetric gauge theory and the Yangian};
Costello-Gwilliam 2017/2021, \emph{Factorization Algebras in Quantum
Field Theory}.

\end{document}
